\documentclass[11pt, a4paper]{article}

% ── Paquetes ──────────────────────────────────────────────────────────────────
\usepackage[utf8]{inputenc}
\usepackage[T1]{fontenc}
\usepackage[spanish]{babel}
\usepackage{geometry}
\geometry{margin=2.5cm}
\usepackage{graphicx}
\usepackage{xcolor}
\usepackage{enumitem}
\usepackage{booktabs}
\usepackage{tabularx}
\usepackage{listings}
\usepackage{amsmath}
\usepackage{hyperref}
\usepackage{fancyhdr}
\usepackage{titlesec}
\usepackage{tcolorbox}
\usepackage{fontawesome5}
\usepackage{tikz}
\usepackage{array}
\usetikzlibrary{arrows.meta, positioning, shapes.geometric}

% ── Colores ───────────────────────────────────────────────────────────────────
\definecolor{primary}{HTML}{1A237E}
\definecolor{accent}{HTML}{E65100}
\definecolor{success}{HTML}{2E7D32}
\definecolor{danger}{HTML}{C62828}
\definecolor{warning}{HTML}{EF6C00}
\definecolor{info}{HTML}{1565C0}
\definecolor{codebg}{HTML}{F5F5F5}
\definecolor{codeframe}{HTML}{BDBDBD}
\definecolor{get}{HTML}{2E7D32}
\definecolor{post}{HTML}{1565C0}
\definecolor{patch}{HTML}{EF6C00}

% ── Configuración de listings ─────────────────────────────────────────────────
\lstset{
  basicstyle=\ttfamily\small,
  backgroundcolor=\color{codebg},
  frame=single,
  rulecolor=\color{codeframe},
  breaklines=true,
  showstringspaces=false,
  extendedchars=true,
  inputencoding=utf8,
  literate=
    {á}{{\'a}}{1} {é}{{\'e}}{1} {í}{{\'i}}{1} {ó}{{\'o}}{1} {ú}{{\'u}}{1}
    {ñ}{{\~n}}{1} {Á}{{\'A}}{1} {É}{{\'E}}{1} {Í}{{\'I}}{1} {Ó}{{\'O}}{1}
    {Ú}{{\'U}}{1} {Ñ}{{\~N}}{1}
}

% ── Cajas (solo para texto, NO para código) ───────────────────────────────────
\newtcolorbox{infobox}[1][]{
  colback=info!5, colframe=info!80!black,
  fonttitle=\bfseries, title={\faInfoCircle\quad Nota}, #1
}

\newtcolorbox{warningbox}[1][]{
  colback=warning!5, colframe=warning!80!black,
  fonttitle=\bfseries, title={\faExclamationCircle\quad Importante}, #1
}

% ── Comando para método HTTP ──────────────────────────────────────────────────
\newcommand{\httpmethod}[2]{%
  \colorbox{#1}{\textcolor{white}{\textbf{\small\textsf{#2}}}}%
}

% ── Comando para encabezado de endpoint ───────────────────────────────────────
\newcommand{\routeheader}[3]{%
  \noindent\rule{\textwidth}{1.5pt}\\[0.3em]
  \httpmethod{#1}{#2}\;\texttt{\bfseries #3}\\[0.3em]
  \rule{\textwidth}{0.4pt}
  \vspace{0.3cm}
}

% ── Comando para encabezado de contrato ───────────────────────────────────────
\newcommand{\contractheader}[1]{%
  \noindent\colorbox{primary!10}{%
    \parbox{\dimexpr\textwidth-2\fboxsep}{%
      \textbf{\faFileCode\quad #1}%
    }%
  }
  \vspace{0.3cm}
}

% ── Encabezado y pie de página ────────────────────────────────────────────────
\pagestyle{fancy}
\fancyhf{}
\fancyhead[L]{\small\textcolor{primary}{\textbf{Grupo 6} --- Contrato API Despacho y Logística}}
\fancyhead[R]{\small\textcolor{primary}{v1.0}}
\fancyfoot[C]{\small\textcolor{gray}{\thepage}}
\renewcommand{\headrulewidth}{0.4pt}

% ── Formato de títulos ────────────────────────────────────────────────────────
\titleformat{\section}
  {\Large\bfseries\color{primary}}{\thesection.}{0.5em}{}[\titlerule]
\titleformat{\subsection}
  {\large\bfseries\color{primary!80}}{\thesubsection}{0.5em}{}
\titleformat{\subsubsection}
  {\normalsize\bfseries\color{primary!60}}{\thesubsubsection}{0.5em}{}

\hypersetup{
  colorlinks=true, linkcolor=primary, urlcolor=info, citecolor=accent,
}

% ══════════════════════════════════════════════════════════════════════════════
\begin{document}

% ── Portada ───────────────────────────────────────────────────────────────────
\begin{titlepage}
\centering
\vspace*{2.5cm}

{\Huge\bfseries\textcolor{primary}{Contrato API}}\\[0.3cm]
{\LARGE\textcolor{primary!70}{Servicio de Despacho y Logística}}\\[0.5cm]
{\Large\textcolor{accent}{\textbf{Grupo 6}}}\\[2cm]

{\large Especificación de Endpoints REST}\\[0.2cm]
{\large Contratos JSON de Eventos}\\[0.2cm]
{\large Matriz de Dependencias}\\[3cm]

\begin{tabular}{rl}
  \textbf{Proyecto:}     & Marketplace en la Nube \\
  \textbf{Versión API:}  & \texttt{v1.0} \\
  \textbf{Base URL:}     & \texttt{https://g6-despacho.onrender.com/api/v1} \\
  \textbf{Fecha:}        & 11 de Junio de 2026 \\
  \textbf{Estado:}       & Borrador para revisión con G5 y G9 \\
\end{tabular}

\vfill

\begin{warningbox}[title={\faExclamationCircle\quad Para revisión conjunta}]
  Este documento debe ser validado con los equipos de \textbf{G5 (Pedidos)}
  y \textbf{G9 (Notificaciones)} antes de comenzar la implementación.
  Cualquier cambio acordado se reflejará en la versión 1.1.
\end{warningbox}

\end{titlepage}

\tableofcontents
\newpage


% ══════════════════════════════════════════════════════════════════════════════
\part{Contrato REST API}
% ══════════════════════════════════════════════════════════════════════════════

\section{Información General}

\begin{table}[h!]
\centering
\renewcommand{\arraystretch}{1.4}
\begin{tabularx}{\textwidth}{l X}
  \toprule
  \textbf{Campo} & \textbf{Valor} \\
  \midrule
  Título          & API de Despacho y Logística --- Grupo 6 \\
  Versión         & 1.0.0 \\
  Base URL        & \texttt{https://g6-despacho.onrender.com/api/v1} \\
  Formato         & JSON (\texttt{application/json}) \\
  Autenticación   & Ninguna (proyecto académico) \\
  Encoding        & UTF-8 \\
  \bottomrule
\end{tabularx}
\end{table}

\subsection{Convenciones}

\begin{itemize}[noitemsep]
  \item Todos los endpoints retornan \texttt{application/json}.
  \item Las fechas usan formato ISO 8601: \texttt{YYYY-MM-DDTHH:mm:ssZ}.
  \item Los IDs de shipment usan formato \texttt{SHP-<uuid>} (generados por nuestro servicio).
  \item Los errores siguen estructura estándar (ver Sección~\ref{sec:errores}).
\end{itemize}

\subsection{Resumen de Endpoints}

\begin{table}[h!]
\centering
\renewcommand{\arraystretch}{1.5}
\begin{tabularx}{\textwidth}{l l X l}
  \toprule
  \textbf{Método} & \textbf{Ruta} & \textbf{Descripción} & \textbf{Consumidor} \\
  \midrule
  \httpmethod{post}{POST}   & \texttt{/shipments}                  & Crear un nuevo despacho          & G5 \\
  \httpmethod{get}{GET}     & \texttt{/shipments/\{id\}}           & Consultar estado de un despacho  & G1, G5 \\
  \httpmethod{get}{GET}     & \texttt{/shipments?order\_id=\{x\}}  & Buscar despacho por pedido       & G5 \\
  \httpmethod{patch}{PATCH} & \texttt{/shipments/\{id\}}           & Actualizar estado del despacho   & Interno \\
  \httpmethod{get}{GET}     & \texttt{/shipments}                  & Listar todos los despachos       & G10 \\
  \httpmethod{get}{GET}     & \texttt{/health}                     & Health check del servicio        & Todos \\
  \bottomrule
\end{tabularx}
\caption{Resumen de endpoints expuestos por el Grupo 6.}
\end{table}


\section{Modelo de Datos: Shipment}
\label{sec:modelo}

\begin{table}[h!]
\centering
\renewcommand{\arraystretch}{1.3}
\begin{tabularx}{\textwidth}{l l l c X}
  \toprule
  \textbf{Campo} & \textbf{Tipo} & \textbf{Ejemplo} & \textbf{Req.} & \textbf{Descripción} \\
  \midrule
  \texttt{shipment\_id}        & string   & \texttt{SHP-a1b2c3d4}        & ---  & ID único generado por G6. \\
  \texttt{order\_id}           & string   & \texttt{ORD-20260611-001}     & Sí   & ID del pedido (de G5). \\
  \texttt{customer\_name}      & string   & \texttt{María G.}             & Sí   & Nombre del destinatario. \\
  \texttt{address}             & string   & \texttt{Av. Prov. 1234}       & Sí   & Dirección de entrega. \\
  \texttt{city}                & string   & \texttt{Santiago}             & Sí   & Ciudad de destino. \\
  \texttt{weight\_kg}          & number   & \texttt{3.2}                  & Sí   & Peso del paquete en kg. \\
  \texttt{status}              & string   & \texttt{PENDING}              & ---  & Estado actual. \\
  \texttt{created\_at}         & datetime & \texttt{2026-06-11T14:00Z}    & ---  & Fecha de creación. \\
  \texttt{updated\_at}         & datetime & \texttt{2026-06-11T15:30Z}    & ---  & Última actualización. \\
  \texttt{estimated\_delivery} & datetime & \texttt{2026-06-14T18:00Z}    & ---  & Fecha estimada. \\
  \bottomrule
\end{tabularx}
\caption{Esquema del recurso \texttt{Shipment}. ``Req.'' indica si es requerido en el \texttt{POST}.}
\end{table}


\subsection{Estados válidos}
\label{sec:estados}

\begin{center}
\begin{tikzpicture}[
  node distance=2.6cm,
  state/.style={rectangle, rounded corners=5pt, draw=primary, thick,
    fill=primary!8, minimum width=2.6cm, minimum height=0.9cm,
    font=\small\bfseries, text=primary},
  error/.style={rectangle, rounded corners=5pt, draw=danger, thick,
    fill=danger!8, minimum width=2.6cm, minimum height=0.9cm,
    font=\small\bfseries, text=danger},
  arrow/.style={-{Stealth[length=3mm]}, thick, primary!70},
  errarrow/.style={-{Stealth[length=3mm]}, thick, danger!70, dashed},
]
  \node[state]                     (pending)   {PENDING};
  \node[state, right of=pending]   (transit)   {IN\_TRANSIT};
  \node[state, right of=transit]   (delivered) {DELIVERED};
  \node[error, below of=pending]   (cancelled) {CANCELLED};
  \node[error, below of=transit]   (failed)    {FAILED};
  \node[error, below of=delivered] (returned)  {RETURNED};
  \draw[arrow]    (pending)   -- (transit);
  \draw[arrow]    (transit)   -- (delivered);
  \draw[errarrow] (pending)   -- (cancelled);
  \draw[errarrow] (transit)   -- (failed);
  \draw[errarrow] (delivered) -- (returned);
\end{tikzpicture}
\end{center}

\begin{table}[h!]
\centering
\renewcommand{\arraystretch}{1.3}
\begin{tabularx}{\textwidth}{l l l X}
  \toprule
  \textbf{Estado} & \textbf{Valor} & \textbf{Tipo} & \textbf{Transiciones válidas} \\
  \midrule
  Pendiente   & \texttt{PENDING}     & Inicial & $\rightarrow$ \texttt{IN\_TRANSIT}, \texttt{CANCELLED} \\
  En tránsito & \texttt{IN\_TRANSIT} & Normal  & $\rightarrow$ \texttt{DELIVERED}, \texttt{FAILED} \\
  Entregado   & \texttt{DELIVERED}   & Final   & $\rightarrow$ \texttt{RETURNED} \\
  Cancelado   & \texttt{CANCELLED}   & Final   & Ninguna (terminal). \\
  Fallido     & \texttt{FAILED}      & Final   & Ninguna (terminal). \\
  Devuelto    & \texttt{RETURNED}    & Final   & Ninguna (terminal). \\
  \bottomrule
\end{tabularx}
\caption{Tabla de transiciones de estado.}
\end{table}


% ══════════════════════════════════════════════════════════════════════════════
\section{Endpoints Detallados}
% ══════════════════════════════════════════════════════════════════════════════

% ── POST /shipments ───────────────────────────────────────────────────────────
\subsection{Crear Despacho}

\routeheader{post}{POST}{/api/v1/shipments}

\textbf{Descripción:} Crea una nueva orden de despacho. Llamado por
\textbf{G5 (Pedidos)} cuando un pedido ha sido confirmado y pagado.

\vspace{0.3cm}
\textbf{Request Body} (\texttt{application/json}):

\begin{lstlisting}
{
  "order_id": "ORD-20260611-001",
  "customer_name": "María González",
  "address": "Av. Providencia 1234, Depto 56",
  "city": "Santiago",
  "weight_kg": 3.2
}
\end{lstlisting}

\textbf{Campos requeridos:}
\begin{table}[h!]
\small
\begin{tabularx}{\textwidth}{l l l X}
  \toprule
  \textbf{Campo} & \textbf{Tipo} & \textbf{Req.} & \textbf{Validaciones} \\
  \midrule
  \texttt{order\_id}      & string & Sí & No vacío. Único. \\
  \texttt{customer\_name} & string & Sí & No vacío. Máx. 200 chars. \\
  \texttt{address}        & string & Sí & No vacío. Máx. 500 chars. \\
  \texttt{city}           & string & Sí & No vacío. Máx. 100 chars. \\
  \texttt{weight\_kg}     & number & Sí & Mayor a 0. Máx. 500 kg. \\
  \bottomrule
\end{tabularx}
\end{table}

\textbf{Response --- 201 Created:}

\begin{lstlisting}
{
  "shipment_id": "SHP-a1b2c3d4",
  "order_id": "ORD-20260611-001",
  "customer_name": "María González",
  "address": "Av. Providencia 1234, Depto 56",
  "city": "Santiago",
  "weight_kg": 3.2,
  "status": "PENDING",
  "created_at": "2026-06-11T14:00:00Z",
  "updated_at": "2026-06-11T14:00:00Z",
  "estimated_delivery": "2026-06-14T18:00:00Z"
}
\end{lstlisting}

\textbf{Errores:}
\begin{table}[h!]
\small
\begin{tabularx}{\textwidth}{l l X}
  \toprule
  \textbf{Código} & \textbf{Nombre} & \textbf{Causa} \\
  \midrule
  400 & Bad Request          & JSON mal formado o campos faltantes. \\
  409 & Conflict             & Ya existe despacho para ese \texttt{order\_id}. \\
  422 & Unprocessable Entity & Datos inválidos (ej: peso negativo). \\
  \bottomrule
\end{tabularx}
\end{table}

\noindent\rule{\textwidth}{0.4pt}

% ── GET /shipments/{id} ──────────────────────────────────────────────────────
\subsection{Consultar Despacho por ID}

\routeheader{get}{GET}{/api/v1/shipments/\{shipment\_id\}}

\textbf{Descripción:} Retorna estado actual y datos completos de un despacho.
Llamado por \textbf{G1 (Frontend)} y \textbf{G5 (Pedidos)}.

\vspace{0.3cm}
\textbf{Response --- 200 OK:}

\begin{lstlisting}
{
  "shipment_id": "SHP-a1b2c3d4",
  "order_id": "ORD-20260611-001",
  "customer_name": "María González",
  "address": "Av. Providencia 1234, Depto 56",
  "city": "Santiago",
  "weight_kg": 3.2,
  "status": "IN_TRANSIT",
  "created_at": "2026-06-11T14:00:00Z",
  "updated_at": "2026-06-11T14:30:00Z",
  "estimated_delivery": "2026-06-14T18:00:00Z"
}
\end{lstlisting}

\textbf{Errores:}
\begin{table}[h!]
\small
\begin{tabularx}{\textwidth}{l l X}
  \toprule
  \textbf{Código} & \textbf{Nombre} & \textbf{Causa} \\
  \midrule
  404 & Not Found & No existe despacho con ese ID. \\
  \bottomrule
\end{tabularx}
\end{table}

\noindent\rule{\textwidth}{0.4pt}

% ── GET /shipments?order_id= ─────────────────────────────────────────────────
\subsection{Buscar Despacho por Order ID}

\routeheader{get}{GET}{/api/v1/shipments?order\_id=\{order\_id\}}

\textbf{Descripción:} Busca un despacho por el ID de pedido de G5.

\vspace{0.3cm}
\textbf{Response --- 200 OK:}

\begin{lstlisting}
{
  "shipment_id": "SHP-a1b2c3d4",
  "order_id": "ORD-20260611-001",
  "status": "PENDING",
  "created_at": "2026-06-11T14:00:00Z",
  "estimated_delivery": "2026-06-14T18:00:00Z"
}
\end{lstlisting}

\textbf{Errores:}
\begin{table}[h!]
\small
\begin{tabularx}{\textwidth}{l l X}
  \toprule
  \textbf{Código} & \textbf{Nombre} & \textbf{Causa} \\
  \midrule
  400 & Bad Request & Falta parámetro \texttt{order\_id}. \\
  404 & Not Found   & No existe despacho para ese pedido. \\
  \bottomrule
\end{tabularx}
\end{table}

\noindent\rule{\textwidth}{0.4pt}

% ── PATCH /shipments/{id} ────────────────────────────────────────────────────
\subsection{Actualizar Estado del Despacho}

\routeheader{patch}{PATCH}{/api/v1/shipments/\{shipment\_id\}}

\textbf{Descripción:} Actualiza el estado de un despacho. Uso \textbf{interno}
(simulación automática) o por \textbf{G5} para cancelar.

\vspace{0.3cm}
\textbf{Request Body:}

\begin{lstlisting}
{
  "status": "IN_TRANSIT"
}
\end{lstlisting}

\textbf{Response --- 200 OK:}

\begin{lstlisting}
{
  "shipment_id": "SHP-a1b2c3d4",
  "order_id": "ORD-20260611-001",
  "status": "IN_TRANSIT",
  "previous_status": "PENDING",
  "updated_at": "2026-06-11T14:30:00Z"
}
\end{lstlisting}

\begin{infobox}
  Al cambiar el estado exitosamente, el servicio emite automáticamente
  el evento correspondiente hacia G9 (ver Parte~II).
\end{infobox}

\textbf{Errores:}
\begin{table}[h!]
\small
\begin{tabularx}{\textwidth}{l l X}
  \toprule
  \textbf{Código} & \textbf{Nombre} & \textbf{Causa} \\
  \midrule
  400 & Bad Request & Estado no es un valor válido. \\
  404 & Not Found   & No existe despacho con ese ID. \\
  409 & Conflict    & Transición no permitida (ej: \texttt{DELIVERED} $\rightarrow$ \texttt{PENDING}). \\
  \bottomrule
\end{tabularx}
\end{table}

\noindent\rule{\textwidth}{0.4pt}

% ── GET /shipments (lista) ───────────────────────────────────────────────────
\subsection{Listar Despachos}

\routeheader{get}{GET}{/api/v1/shipments}

\textbf{Descripción:} Lista de despachos con filtros. Para \textbf{G10 (Reportería)}.

\vspace{0.3cm}
\textbf{Query Parameters (opcionales):}
\begin{table}[h!]
\small
\begin{tabularx}{\textwidth}{l l l X}
  \toprule
  \textbf{Param} & \textbf{Tipo} & \textbf{Default} & \textbf{Descripción} \\
  \midrule
  \texttt{status} & string & (todos) & Filtrar por estado. \\
  \texttt{limit}  & int    & 50      & Máx. resultados. \\
  \texttt{offset} & int    & 0       & Paginación. \\
  \bottomrule
\end{tabularx}
\end{table}

\textbf{Response --- 200 OK:}

\begin{lstlisting}
{
  "total": 142,
  "limit": 50,
  "offset": 0,
  "shipments": [
    {
      "shipment_id": "SHP-a1b2c3d4",
      "order_id": "ORD-20260611-001",
      "status": "DELIVERED",
      "city": "Santiago",
      "created_at": "2026-06-11T14:00:00Z",
      "updated_at": "2026-06-11T16:00:00Z"
    }
  ]
}
\end{lstlisting}

\noindent\rule{\textwidth}{0.4pt}

% ── GET /health ──────────────────────────────────────────────────────────────
\subsection{Health Check}

\routeheader{get}{GET}{/api/v1/health}

\begin{lstlisting}
{
  "status": "ok",
  "service": "g6-despacho",
  "version": "1.0.0",
  "timestamp": "2026-06-11T14:00:00Z"
}
\end{lstlisting}

\noindent\rule{\textwidth}{0.4pt}


% ── Estructura Estándar de Errores ───────────────────────────────────────────
\section{Estructura Estándar de Errores}
\label{sec:errores}

Todos los errores siguen el mismo formato:

\begin{lstlisting}
{
  "error": {
    "code": 422,
    "type": "VALIDATION_ERROR",
    "message": "El campo weight_kg debe ser mayor a 0.",
    "details": {
      "field": "weight_kg",
      "value": -1.5,
      "constraint": "must be greater than 0"
    }
  }
}
\end{lstlisting}

\begin{table}[h!]
\centering
\renewcommand{\arraystretch}{1.3}
\begin{tabularx}{\textwidth}{l l X}
  \toprule
  \textbf{Código} & \textbf{type} & \textbf{Uso} \\
  \midrule
  400 & \texttt{BAD\_REQUEST}      & JSON malformado, campos faltantes. \\
  404 & \texttt{NOT\_FOUND}        & Recurso no encontrado. \\
  409 & \texttt{CONFLICT}          & Duplicado o transición inválida. \\
  422 & \texttt{VALIDATION\_ERROR} & Datos que no pasan validación. \\
  500 & \texttt{INTERNAL\_ERROR}   & Error inesperado del servidor. \\
  \bottomrule
\end{tabularx}
\caption{Códigos de error estandarizados.}
\end{table}


% ══════════════════════════════════════════════════════════════════════════════
\newpage
\part{Contrato de Eventos JSON}
% ══════════════════════════════════════════════════════════════════════════════

\section{Mecanismo de Comunicación}

\begin{warningbox}[title={\faExclamationCircle\quad Decisión pendiente con G9}]
  Se propone usar \textbf{Webhooks HTTP}: nuestro servicio hace un
  \texttt{POST} a una URL de G9 cada vez que cambia un estado.

  \vspace{0.3em}
  \textbf{Alternativa:} G9 hace polling a nuestro \texttt{GET}.

  \vspace{0.3em}
  \textbf{Acordar en la reunión de contratos antes del 16 de Junio.}
\end{warningbox}

\subsection{Flujo de Webhooks (propuesto)}

\begin{enumerate}[noitemsep]
  \item G9 nos registra su URL de webhook.
  \item Cada cambio de estado dispara un \texttt{POST} a esa URL.
  \item G9 responde \texttt{200 OK} para confirmar recepción.
  \item Si falla, reintentamos 3 veces con backoff (1s, 2s, 4s).
\end{enumerate}


\section{Formato Estándar del Evento}
\label{sec:formato-evento}

Todos los eventos emitidos por G6 siguen esta estructura:

\contractheader{Envelope del Evento (estructura común)}

\begin{lstlisting}
{
  "event_id": "evt-550e8400-e29b-41d4-a716-446655440000",
  "event_type": "ShipmentDelivered",
  "source": "g6-despacho",
  "timestamp": "2026-06-11T15:00:00Z",
  "payload": { ... }
}
\end{lstlisting}

\begin{table}[h!]
\centering
\renewcommand{\arraystretch}{1.3}
\begin{tabularx}{\textwidth}{l l X}
  \toprule
  \textbf{Campo} & \textbf{Tipo} & \textbf{Descripción} \\
  \midrule
  \texttt{event\_id}   & string (UUID) & Identificador único del evento. \\
  \texttt{event\_type} & string        & Tipo del evento (ver catálogo). \\
  \texttt{source}      & string        & Siempre \texttt{"g6-despacho"}. \\
  \texttt{timestamp}   & datetime      & Momento del evento (ISO 8601). \\
  \texttt{payload}     & object        & Datos específicos. \\
  \bottomrule
\end{tabularx}
\caption{Campos del envelope de evento.}
\end{table}


\section{Catálogo de Eventos}

% ── ShipmentCreated ──────────────────────────────────────────────────────────
\subsection{ShipmentCreated}

Se emite cuando G5 crea exitosamente un nuevo despacho.

\contractheader{ShipmentCreated}

\begin{lstlisting}
{
  "event_id": "evt-550e8400-e29b-41d4-a716-446655440001",
  "event_type": "ShipmentCreated",
  "source": "g6-despacho",
  "timestamp": "2026-06-11T14:00:00Z",
  "payload": {
    "shipment_id": "SHP-a1b2c3d4",
    "order_id": "ORD-20260611-001",
    "customer_name": "María González",
    "address": "Av. Providencia 1234, Depto 56",
    "city": "Santiago",
    "weight_kg": 3.2,
    "status": "PENDING",
    "estimated_delivery": "2026-06-14T18:00:00Z"
  }
}
\end{lstlisting}

\textbf{Consumidores:} G9 (notifica al cliente), G10 (registro).

\vspace{0.4cm}

% ── ShipmentInTransit ────────────────────────────────────────────────────────
\subsection{ShipmentInTransit}

Se emite cuando el paquete sale a reparto.

\contractheader{ShipmentInTransit}

\begin{lstlisting}
{
  "event_id": "evt-550e8400-e29b-41d4-a716-446655440002",
  "event_type": "ShipmentInTransit",
  "source": "g6-despacho",
  "timestamp": "2026-06-11T14:30:00Z",
  "payload": {
    "shipment_id": "SHP-a1b2c3d4",
    "order_id": "ORD-20260611-001",
    "previous_status": "PENDING",
    "new_status": "IN_TRANSIT",
    "customer_name": "María González",
    "city": "Santiago"
  }
}
\end{lstlisting}

\textbf{Consumidores:} G9 (``Tu paquete va en camino'').

\vspace{0.4cm}

% ── ShipmentDelivered ────────────────────────────────────────────────────────
\subsection{ShipmentDelivered}

Se emite cuando el paquete es entregado exitosamente.

\contractheader{ShipmentDelivered}

\begin{lstlisting}
{
  "event_id": "evt-550e8400-e29b-41d4-a716-446655440003",
  "event_type": "ShipmentDelivered",
  "source": "g6-despacho",
  "timestamp": "2026-06-11T15:00:00Z",
  "payload": {
    "shipment_id": "SHP-a1b2c3d4",
    "order_id": "ORD-20260611-001",
    "previous_status": "IN_TRANSIT",
    "new_status": "DELIVERED",
    "customer_name": "María González",
    "city": "Santiago",
    "delivered_at": "2026-06-11T15:00:00Z"
  }
}
\end{lstlisting}

\textbf{Consumidores:} G9 (``Tu paquete ha llegado''), G10 (métricas).

\vspace{0.4cm}

% ── ShipmentCancelled ────────────────────────────────────────────────────────
\subsection{ShipmentCancelled}

Se emite cuando un despacho es cancelado antes de salir a reparto.

\contractheader{ShipmentCancelled}

\begin{lstlisting}
{
  "event_id": "evt-550e8400-e29b-41d4-a716-446655440004",
  "event_type": "ShipmentCancelled",
  "source": "g6-despacho",
  "timestamp": "2026-06-11T14:10:00Z",
  "payload": {
    "shipment_id": "SHP-a1b2c3d4",
    "order_id": "ORD-20260611-001",
    "previous_status": "PENDING",
    "new_status": "CANCELLED",
    "customer_name": "María González",
    "reason": "Pedido anulado por el cliente"
  }
}
\end{lstlisting}

\textbf{Consumidores:} G9 (notifica cancelación), G5 (confirma).

\vspace{0.4cm}

% ── ShipmentFailed ───────────────────────────────────────────────────────────
\subsection{ShipmentFailed}

Se emite cuando un intento de entrega falla.

\contractheader{ShipmentFailed}

\begin{lstlisting}
{
  "event_id": "evt-550e8400-e29b-41d4-a716-446655440005",
  "event_type": "ShipmentFailed",
  "source": "g6-despacho",
  "timestamp": "2026-06-11T16:00:00Z",
  "payload": {
    "shipment_id": "SHP-a1b2c3d4",
    "order_id": "ORD-20260611-001",
    "previous_status": "IN_TRANSIT",
    "new_status": "FAILED",
    "customer_name": "María González",
    "reason": "Destinatario ausente"
  }
}
\end{lstlisting}

\textbf{Consumidores:} G9 (``Hubo un problema con tu entrega''), G5.

\vspace{0.4cm}

% ── ShipmentReturned ─────────────────────────────────────────────────────────
\subsection{ShipmentReturned}

Se emite cuando un paquete entregado es devuelto.

\contractheader{ShipmentReturned}

\begin{lstlisting}
{
  "event_id": "evt-550e8400-e29b-41d4-a716-446655440006",
  "event_type": "ShipmentReturned",
  "source": "g6-despacho",
  "timestamp": "2026-06-12T10:00:00Z",
  "payload": {
    "shipment_id": "SHP-a1b2c3d4",
    "order_id": "ORD-20260611-001",
    "previous_status": "DELIVERED",
    "new_status": "RETURNED",
    "customer_name": "María González",
    "reason": "Producto incorrecto"
  }
}
\end{lstlisting}

\textbf{Consumidores:} G9, G5 (gestionar reembolso), G10 (métricas).


% ══════════════════════════════════════════════════════════════════════════════
\newpage
\part{Contratos Inter-Grupo}
% ══════════════════════════════════════════════════════════════════════════════

\section{Contrato con G5 (Pedidos) --- Entrada}
\label{sec:contrato-g5}

Este contrato define \textbf{qué datos nos envía G5} cuando un pedido es
confirmado y necesita ser despachado.

\subsection{Flujo de Integración}

\begin{enumerate}[noitemsep]
  \item El cliente completa el pago en G5.
  \item G5 hace una solicitud HTTP POST a nuestro \texttt{/api/v1/shipments}.
  \item Respondemos con el \texttt{shipment\_id} y estado \texttt{PENDING}.
  \item G5 almacena el \texttt{shipment\_id} para consultas futuras.
\end{enumerate}

\subsection{JSON que G5 nos envía}

\contractheader{Request de G5 hacia G6}

\begin{lstlisting}
{
  "order_id": "ORD-20260611-001",
  "customer_name": "María González",
  "address": "Av. Providencia 1234, Depto 56",
  "city": "Santiago",
  "weight_kg": 3.2
}
\end{lstlisting}

\subsection{JSON que G6 le responde a G5}

\contractheader{Response de G6 hacia G5}

\begin{lstlisting}
{
  "shipment_id": "SHP-a1b2c3d4",
  "order_id": "ORD-20260611-001",
  "customer_name": "María González",
  "address": "Av. Providencia 1234, Depto 56",
  "city": "Santiago",
  "weight_kg": 3.2,
  "status": "PENDING",
  "created_at": "2026-06-11T14:00:00Z",
  "updated_at": "2026-06-11T14:00:00Z",
  "estimated_delivery": "2026-06-14T18:00:00Z"
}
\end{lstlisting}

\subsection{Preguntas pendientes para G5}

\begin{warningbox}[title={\faExclamationCircle\quad Confirmar con G5}]
\begin{enumerate}[noitemsep]
  \item Formato de \texttt{order\_id}: propuesto como string libre.
  \item ¿Necesitan campos adicionales en la respuesta?
  \item ¿Nos envían el peso o lo calculamos nosotros?
  \item ¿Cancelan vía \texttt{PATCH} o necesitan endpoint dedicado?
\end{enumerate}
\end{warningbox}


\section{Contrato con G9 (Notificaciones) --- Salida}
\label{sec:contrato-g9}

Define \textbf{qué datos le enviamos a G9} en cada cambio de estado.

\subsection{Flujo de Integración}

\begin{enumerate}[noitemsep]
  \item Un shipment cambia de estado (vía PATCH o simulación).
  \item Construimos el evento JSON correspondiente.
  \item Hacemos \texttt{POST} al webhook de G9.
  \item G9 procesa y notifica al usuario final.
\end{enumerate}

\subsection{JSON que G6 le envía a G9}

\contractheader{Evento de G6 hacia G9 (ejemplo: entrega)}

\begin{lstlisting}
{
  "event_id": "evt-550e8400-e29b-41d4-a716-446655440003",
  "event_type": "ShipmentDelivered",
  "source": "g6-despacho",
  "timestamp": "2026-06-11T15:00:00Z",
  "payload": {
    "shipment_id": "SHP-a1b2c3d4",
    "order_id": "ORD-20260611-001",
    "previous_status": "IN_TRANSIT",
    "new_status": "DELIVERED",
    "customer_name": "María González",
    "city": "Santiago",
    "delivered_at": "2026-06-11T15:00:00Z"
  }
}
\end{lstlisting}

\subsection{Respuesta esperada de G9}

\contractheader{Acknowledgment de G9}

\begin{lstlisting}
{
  "received": true,
  "event_id": "evt-550e8400-e29b-41d4-a716-446655440003"
}
\end{lstlisting}

\subsection{Preguntas pendientes para G9}

\begin{warningbox}[title={\faExclamationCircle\quad Confirmar con G9}]
\begin{enumerate}[noitemsep]
  \item ¿URL de webhook?
  \item ¿Campos adicionales necesarios? (email, teléfono)
  \item ¿Prefieren webhooks o polling?
  \item ¿Cómo confirman recepción? Proponemos \texttt{200 OK}.
  \item ¿Qué eventos les interesan? ¿Los 6 o solo algunos?
\end{enumerate}
\end{warningbox}


\section{Contrato con G10 (Reportería) --- Salida}

G10 consume datos para dashboards de rendimiento.

\subsection{Opciones de Integración}

\begin{enumerate}[noitemsep]
  \item \textbf{Vía API REST:} \texttt{GET /shipments} con filtros.
  \item \textbf{Vía Eventos:} Suscripción a webhooks (mismos que G9).
\end{enumerate}

\subsection{Datos relevantes para G10}

\begin{table}[h!]
\centering
\renewcommand{\arraystretch}{1.3}
\begin{tabularx}{\textwidth}{l X}
  \toprule
  \textbf{Métrica} & \textbf{Cómo la obtienen} \\
  \midrule
  Total despachos       & \texttt{GET /shipments} $\rightarrow$ \texttt{total}. \\
  Por estado            & Filtrar con \texttt{?status=DELIVERED}. \\
  Tasa de éxito         & \texttt{DELIVERED} / total. \\
  Tasa de fallos        & (\texttt{FAILED} + \texttt{RETURNED}) / total. \\
  Tiempo de entrega     & \texttt{updated\_at} - \texttt{created\_at} en \texttt{DELIVERED}. \\
  \bottomrule
\end{tabularx}
\end{table}


\section{Contrato con G1 (Frontend) --- Consulta}

G1 consume nuestro endpoint de tracking.

\textbf{Endpoint:} \texttt{GET /api/v1/shipments/\{shipment\_id\}}

\begin{table}[h!]
\centering
\renewcommand{\arraystretch}{1.3}
\begin{tabularx}{\textwidth}{l X}
  \toprule
  \textbf{Campo} & \textbf{Uso en la UI} \\
  \midrule
  \texttt{status}              & Paso actual en la barra de progreso. \\
  \texttt{estimated\_delivery} & ``Llega el 14 de Junio''. \\
  \texttt{city}                & ``Tu paquete va camino a Santiago''. \\
  \texttt{updated\_at}         & ``Última actualización: hace 30 min''. \\
  \bottomrule
\end{tabularx}
\end{table}


% ══════════════════════════════════════════════════════════════════════════════
\newpage
\part{Matriz de Dependencias y Arquitectura}
% ══════════════════════════════════════════════════════════════════════════════

\section{Matriz de Dependencias}

\begin{table}[h!]
\centering
\renewcommand{\arraystretch}{1.5}
\begin{tabularx}{\textwidth}{c c l l X}
  \toprule
  \textbf{Origen} & \textbf{Destino} & \textbf{Tipo} & \textbf{Protocolo} & \textbf{Datos} \\
  \midrule
  G5  & G6  & Síncrono  & REST POST    & Datos del pedido. \\
  G5  & G6  & Síncrono  & REST GET     & Consulta por \texttt{order\_id}. \\
  G5  & G6  & Síncrono  & REST PATCH   & Cancelación. \\
  G1  & G6  & Síncrono  & REST GET     & Tracking por \texttt{shipment\_id}. \\
  G6  & G9  & Asíncrono & Webhook POST & Eventos de estado. \\
  G6  & G10 & Asíncrono & Webhook POST & Eventos de estado. \\
  G10 & G6  & Síncrono  & REST GET     & Lista para reportes. \\
  \bottomrule
\end{tabularx}
\caption{Matriz completa de dependencias del Grupo 6.}
\end{table}


\section{Diagrama de Arquitectura}

\begin{center}
\begin{tikzpicture}[
  node distance=2.5cm and 3.5cm,
  box/.style={rectangle, rounded corners=8pt, draw=primary, thick, fill=primary!6,
    minimum width=3cm, minimum height=1.4cm, align=center,
    font=\small\bfseries, text=primary},
  extbox/.style={rectangle, rounded corners=8pt, draw=info, thick, fill=info!6,
    minimum width=3cm, minimum height=1.4cm, align=center,
    font=\small\bfseries, text=info},
  db/.style={cylinder, draw=accent, thick, fill=accent!8, shape border rotate=90,
    minimum width=2.2cm, minimum height=1.2cm, aspect=0.3,
    font=\small\bfseries, text=accent},
  arr/.style={-{Stealth[length=3mm]}, thick},
]
  \node[extbox]                    (g1)  {G1\\\relax Frontend};
  \node[extbox, below=1.5cm of g1] (g5)  {G5\\\relax Pedidos};
  \node[box, right=3cm of g5]      (g6)  {G6\\\relax Despachos};
  \node[db, below=1.5cm of g6]     (db)  {Supabase};
  \node[extbox, right=3cm of g6]   (g9)  {G9\\\relax Notificaciones};
  \node[extbox, below=1.5cm of g9] (g10) {G10\\\relax Reportería};

  \draw[arr, info] (g5)  -- node[above, font=\scriptsize] {\texttt{POST/GET/PATCH}} (g6);
  \draw[arr, info] (g1)  -| node[above left, font=\scriptsize] {\texttt{GET}} (g6);
  \draw[arr, primary] (g6) -- node[left, font=\scriptsize] {SQL} (db);
  \draw[arr, success] (g6) -- node[above, font=\scriptsize] {Webhook} (g9);
  \draw[arr, success] (g6) |- node[below right, font=\scriptsize] {Webhook} (g10);
  \draw[arr, info, dashed] (g10) -- node[right, font=\scriptsize] {\texttt{GET}} (g6);
\end{tikzpicture}
\end{center}


\section{Esquema SQL de la Base de Datos}

\begin{lstlisting}
CREATE TABLE shipments (
    shipment_id        TEXT PRIMARY KEY,
    order_id           TEXT UNIQUE NOT NULL,
    customer_name      TEXT NOT NULL,
    address            TEXT NOT NULL,
    city               TEXT NOT NULL,
    weight_kg          REAL NOT NULL CHECK (weight_kg > 0),
    status             TEXT NOT NULL DEFAULT 'PENDING'
                       CHECK (status IN (
                           'PENDING', 'IN_TRANSIT', 'DELIVERED',
                           'CANCELLED', 'FAILED', 'RETURNED'
                       )),
    created_at         TIMESTAMP NOT NULL DEFAULT NOW(),
    updated_at         TIMESTAMP NOT NULL DEFAULT NOW(),
    estimated_delivery TIMESTAMP
);
\end{lstlisting}

\begin{infobox}
  Esquema para Supabase (PostgreSQL). Si se usa Firebase,
  se adaptará a documentos JSON con los mismos campos.
\end{infobox}


\section{Historial de Versiones}

\begin{table}[h!]
\centering
\begin{tabularx}{\textwidth}{l l l X}
  \toprule
  \textbf{Ver.} & \textbf{Fecha} & \textbf{Autor} & \textbf{Cambios} \\
  \midrule
  1.0 & 11 Jun 2026 & Grupo 6 & Versión inicial. \\
  \bottomrule
\end{tabularx}
\end{table}


\end{document}
